%%%%%%%%%%%%%%%%%%%%%%%%%%%%%%%%%%%%%%%%%%%%%%%%%%%%%%%%%%%%%
% ==================== 5.1：问题一的模型建立与求解 ====================
% ✓ 自查：模型选择有依据；模型已汇总；步骤可复现；结果真实；检验与模型匹配。

\subsection{问题一的模型建立与求解}

\subsubsection{建模思路与模型选择}

问题一要求【概括问题一的任务与目标】，其核心是回答【需要解决的关键数学问题】。
根据【题目条件、研究对象或数据特征】，该问题属于【数学问题类型】。
综合比较【候选方法】在【适用条件、解释能力或求解效率】方面的特点，本文选择
【主模型名称】刻画【变量关系或现实机理】，并以【必要的对比方法】作为参照。
若本问确有数据处理需求，则采用【处理方法】处理【缺失、异常、量纲或口径问题】，理由是【选择依据】；
所得【参数、特征或中间结果】将作为问题二的输入。

\subsubsection{模型的建立}

定义【决策变量、状态变量或解释变量】为【符号及含义】，定义【关键参数】为【符号及含义】。
依据【理论、机理或题目规则】建立【核心方程、目标函数或判别规则】，并结合
【资源限制、边界条件或逻辑关系】给出相应约束：

\begin{equation}
  \text{【核心方程、目标函数或判别规则】},
  \label{eq:q1_core}
\end{equation}

\begin{equation}
  \text{【约束条件、边界条件或适用条件】}.
  \label{eq:q1_constraint}
\end{equation}

\noindent\textbf{模型汇总：}
将上述变量关系、约束条件和参数范围统一整理，得到问题一的完整模型

\begin{equation}
  \mathcal{M}_1:\quad
  \left\{
  \begin{aligned}
    &\text{【核心模型表达式】},\\
    &\text{【约束或适用条件】},\\
    &\text{【变量范围、初始条件或边界条件】}.
  \end{aligned}
  \right.
  \label{eq:q1_summary}
\end{equation}

\subsubsection{模型的求解}

\noindent\textbf{Step 1：}整理【原始输入、题设参数或初始条件】，得到【可直接用于求解的输入形式】。

\noindent\textbf{Step 2：}按照【算法或计算方法】计算【关键变量、中间量或候选解】，并依据【规则】完成【核心求解过程】。

\noindent\textbf{Step 3：}根据【终止条件、筛选准则或评价指标】确定最终结果，并保存【问题二需要调用的输出】。

求解使用【软件、求解器或程序语言】完成，关键参数取值及其来源为【参数设置与依据】，完整程序见附录【编号】。

\subsubsection{求解结果与分析}

求解得到【真实的核心结果】。由【图或表编号】可知，【描述主要现象或变化规律】；
其中【关键结果】说明【数学含义或现实含义】。结合题目要求，问题一的结论为【直接回答问题一】，
并将【真实的中间结果】传递至问题二。

\subsubsection{模型检验与灵敏度分析}

针对【本问模型类型与主要结论】，选取【检验方法】进行验证。
检验指标【指标名称】的结果为【真实结果】，与【判定标准或基准方法】对比后表明【检验结论】。
若【关键参数】具有现实不确定性且会影响主要结论，则在【有依据的变化范围】内进行扰动，
分析【输出数值或决策方案】的变化并说明【稳定区间、临界点或敏感方向】；若无必要，可删除本句而不强行开展灵敏度分析。
